\section{从惊讶到 Bit 与自信息}
\label{sec:07-Information-Theory/01-Information-and-Entropy/01}

同事在纸条上写了一个字母藏起来，让你猜。候选是 \(a\) 到 \(h\) 这八个，他是随机挑的，每个机会均等。规则：你只能问是非题，他只答``是''或``否''。

一个一个试当然可以，运气差要问七次。更好的问法是每次砍掉一半：``在前四个里吗''——无论答什么都只剩四个；再问一次剩两个；第三次锁定。八个候选，三个问题，不多不少。十六个候选就是四个问题。这个数正是候选数的以二为底的对数。

现在把游戏改一改。同事承认他偏心：\(a\) 他会以 \(1/2\) 的概率写下，\(b\) 是 \(1/4\)，\(c\) 和 \(d\) 各 \(1/8\)，其余四个字母他从不写。等分候选的策略在这里就不划算了——第一问应该单独问``是 \(a\) 吗''，因为一半的时候你一问就中。

\subsection{把一次询问画成一棵树}

按照``先问最可能的''这条思路把提问过程展开，得到一棵二叉树。每个字母落在某片叶子上，从根走到叶子的边数，就是猜中它所需的提问次数。

\begin{center}
\begin{tikzpicture}[x=1cm,y=1cm,font=\small]
  \draw[black!45] (0,3.3) -- (-2.3,2.2);
  \draw[black!45] (0,3.3) -- (0.9,2.2);
  \draw[black!45] (0.9,2.2) -- (-0.4,1.1);
  \draw[black!45] (0.9,2.2) -- (2.2,1.1);
  \draw[black!45] (2.2,1.1) -- (1.4,0);
  \draw[black!45] (2.2,1.1) -- (3.0,0);
  \fill[black!70] (0,3.3) circle (2.4pt);
  \fill[black!70] (0.9,2.2) circle (2.4pt);
  \fill[black!70] (2.2,1.1) circle (2.4pt);
  \node[right,font=\footnotesize,black!75] at (0.12,3.48) {是 \(a\) 吗};
  \node[right,font=\footnotesize,black!75] at (1.02,2.38) {是 \(b\) 吗};
  \node[right,font=\footnotesize,black!75] at (2.32,1.28) {是 \(c\) 吗};
  \node[font=\footnotesize,black!55] at (-1.45,2.98) {是};
  \node[font=\footnotesize,black!55] at (0.62,2.92) {否};
  \node[font=\footnotesize,black!55] at (-0.02,1.88) {是};
  \node[font=\footnotesize,black!55] at (1.78,1.86) {否};
  \node[font=\footnotesize,black!55] at (1.56,0.76) {是};
  \node[font=\footnotesize,black!55] at (2.92,0.76) {否};
  \draw[blue!60,fill=blue!8] (-2.3,2.2) circle (0.3);
  \node at (-2.3,2.2) {\(a\)};
  \draw[blue!60,fill=blue!8] (-0.4,1.1) circle (0.3);
  \node at (-0.4,1.1) {\(b\)};
  \draw[blue!60,fill=blue!8] (1.4,0) circle (0.3);
  \node at (1.4,0) {\(c\)};
  \draw[blue!60,fill=blue!8] (3.0,0) circle (0.3);
  \node at (3.0,0) {\(d\)};
  \node[font=\footnotesize,black!70,align=center] at (-2.3,1.55) {\(1/2\)\\问 1 次};
  \node[font=\footnotesize,black!70,align=center] at (-0.4,0.45) {\(1/4\)\\问 2 次};
  \node[font=\footnotesize,black!70,align=center] at (1.4,-0.65) {\(1/8\)\\问 3 次};
  \node[font=\footnotesize,black!70,align=center] at (3.0,-0.65) {\(1/8\)\\问 3 次};
  \draw[black!25] (4.15,-0.9) -- (4.15,3.6);
  \node[right,font=\footnotesize,black!75,align=left] at (4.35,1.9)
    {平均提问次数\\[2pt]\(\tfrac12\cdot1+\tfrac14\cdot2+\tfrac18\cdot3+\tfrac18\cdot3\)\\[2pt]\(=1.75\)};
\end{tikzpicture}

\emph{每片叶子的深度恰好等于 \(-\log_2 p\)：概率 \(1/2\) 的 \(a\) 在第一层，概率 \(1/8\) 的 \(c\) 和 \(d\) 在第三层。把深度按概率加权平均，得到 \(1.75\)——这棵树的平均深度。整个单元都在围绕这张图打转，后面的熵、编码长度、困惑度，说的都是它。}
\end{center}

树上有件事值得盯住：叶子的深度不是随便安排的，它被概率钉死了。\(a\) 占了一半的概率，就配得上最短的路径；\(c\) 只占八分之一，多走两步也不冤枉。写成公式，结果 \(x\) 的定位代价是

\[
\imath_X(x)=-\log_b P(X=x).
\]

这个量叫自信息。观察结果之前它是随机的，观察之后才落成一个具体的数。名字里的``信息''容易让人往语义上想，但它衡量的只有一件事：在给定的概率模型下，把这个结果从所有可能性中拣出来要花多少比特。

\subsection{为什么非得是对数}

换个角度也能到达同一个公式。假设我们只对信息量提三条要求：概率越小信息越大；必然事件的信息为零；两个独立事件同时发生时，各自的信息相加。

第三条最有约束力。独立意味着 \(P(A\cap B)=P(A)P(B)\)，而我们要求 \(I(pq)=I(p)+I(q)\)。把乘法搬成加法的连续函数只有对数的常数倍，这是 Cauchy 函数方程的标准结论。第一条要求把常数的符号固定为负，剩下的自由度被底数 \(b\) 吸收，于是 \(I(p)=-\log_b p\)。

也就是说，对数不是从一堆候选公式里挑出来的，是三条要求把它逼出来的。而提问树给了同一个公式一个操作性的读法：\(-\log_2 p\) 就是最优提问下的深度。两条路走到一起，这个定义才算落了地。

\subsection{Bit 是单位，不是二进制位}

底数决定单位。取 \(2\) 得到 bit，取自然底 \(e\) 得到 nat，取 \(10\) 得到 hartley。换底只是乘一个常数：

\[
\log_b p=\frac{\log_c p}{\log_c b}.
\]

理论推导偏爱 nat，因为微分干净；工程报告偏爱 bit，因为可以直接读成``几个是非题''或``几位存储''。同一个分布用两种单位报出来的数不同，判断不该因此改变——看到一个熵值先问单位，和看到一个长度先问是米还是英尺一样自然。

有个混淆值得先拆掉。一个结果的自信息可以是 \(0.47\) 这样的非整数，而磁盘上的一位就是一位，没有半位。概率为 \(0.9\) 的结果自信息约 \(0.152\) bit，你不可能只写下十分之一个二进制位。矛盾的化解方式是不逐个符号编码：把一长串结果打包，让整串共用一段码，分数位就在平摊中被消化掉了。提问树上对应的操作是一次性问一串字母的联合信息，而不是一个一个问。这条路线在\hyperref[sec:07-Information-Theory/04-Source-Coding/01]{``前缀码与 Kraft 不等式''}里正式展开。

\subsection{独立让惊讶度相加}

若 \(X\) 与 \(Y\) 独立，联合结果的自信息直接拆开：

\[
\imath_{X,Y}(x,y)=-\log\bigl(P(X=x)P(Y=y)\bigr)=\imath_X(x)+\imath_Y(y).
\]

回到提问树，这句话的意思是先问清 \(X\) 再问清 \(Y\)，两段提问互不省力。若两者相关，情况就变了：知道了 \(X\) 之后关于 \(Y\) 的提问会变少，联合惊讶度小于两项之和，而那个差额正是互信息，见\hyperref[sec:07-Information-Theory/02-Joint-Information/02]{``互信息是依赖，也是对不确定性的减少''}。

\subsection{惊讶度不等于重要性，也不等于事实}

自信息只看模型给出的概率，语义、效用、后果一概不管。一段随机噪声难以预测，自信息很高，却没有任何内容；一条``该患者对青霉素过敏''的记录在人群里概率不低，自信息平平，但它决定用药方案。信息论有意把语义扔在门外，换来的是一把在任何概率模型上都能用的尺子——它量不出的恰恰是价值。数据工作里最常见的误用，就是把``熵高''当作``内容更丰富''来汇报。

另一重依赖更隐蔽。写 \(\imath_X(x)\) 时，模型 \(P\) 是隐含参数。如果你的模型认为某事件概率是 \(10^{-6}\)，而它每周都发生，算出来的二十 bit 惊讶度说明的是模型被数据反驳得很惨，不是这件事在真实机制里有多稀有。所以每个自信息值都该读成``相对于分布 \(P\) 的信息''；分布换了，数就换了。用错模型的系统性后果，会在\hyperref[sec:07-Information-Theory/03-Divergences/01]{``错误模型的代价：从交叉熵到 KL 散度''}里被量化成一个具体的额外比特数。

到这里手上只有单个结果的自信息，而且要等结果出现才能算。可实际做决定时，往往在拆开纸条之前就要估计这场游戏值不值得玩。下一节对自信息取期望，把一次询问的平均深度算出来。
